% GENERATED FILE. Edit canonical lessons, then run npm run generate:books.
% canonical-source-hash: fnv1a64:48b09bdee2b64537
% canonical-lessons: BN-C173-jilipi, BN-C173-laddu, BN-C173-gur, BN-C173-jam, BN-C173-lichu

\chapter{Jalebi, Laddu, Jaggery, Black plum, Lychee}
\label{ch:bn-jalebi-laddu-jaggery-black-plum-lychee}
\hlchaptermodality{173}

\begin{chapteropening}
\textbf{By the end of this chapter:} \emph{I can say a jalebi, a laddu, jaggery, a black plum and a lychee.}

Complete the last lesson of chapter 173: I can say a jalebi, a laddu, jaggery, a black plum and a lychee.
\end{chapteropening}

\section[\texorpdfstring{\bn{জিলিপি} (jilipi)}{জিলিপি (jilipi)}]{\bn{জিলিপি} (jilipi) — a jalebi}

\label{lesson:BN-C173-jilipi}



\begin{quote}
Before the new one: say the Bengali for chalk, then the Bengali for a scholarship.
\end{quote}

\subsection*{You'll want to know: \bn{জিলিপি}}
\textbf{\bn{জিলিপি}} — \emph{jilipi} — \textquotedblleft{}a jalebi\textquotedblright{}.

\begin{cousinweb}
\textbf{\bn{জিলিপি}} (\emph{jilipi}) is usually traced, through Persian, to Arabic \emph{zalābiya}, the name of a fried sweet.
\end{cousinweb}

\begin{grammarlens}[title={What you've built}]
One more word for this chapter.
\end{grammarlens}

\subsection*{Your turn}
\begin{itemize}
\raggedright
  \item \emph{Say it:} \emph{jilipi}
  \item \emph{Say it:} \emph{jilipi}, once more
  \item \emph{Say it:} say \emph{jilipi}
  \item \emph{Recall:} say the Bengali for chalk, then the Bengali for a scholarship, then say \emph{jilipi} again
\end{itemize}

\begin{tcolorbox}[breakable,colback=teal!4,colframe=teal!35!black,title={Before you move on}]
What does \bn{জিলিপি} mean? (\textquotedblleft{}A jalebi\textquotedblright{}.) Say it once more.
\end{tcolorbox}
\section[\texorpdfstring{\bn{লাড্ডু} (lāḍḍu)}{লাড্ডু (lāḍḍu)}]{\bn{লাড্ডু} (lāḍḍu) — a laddu}

\label{lesson:BN-C173-laddu}



\begin{quote}
Before the new one: say the Bengali for a scholarship, then the Bengali for a jalebi.
\end{quote}

\subsection*{You'll want to know: \bn{লাড্ডু}}
\textbf{\bn{লাড্ডু}} — \emph{lāḍḍu} — \textquotedblleft{}a laddu\textquotedblright{}.

\begin{cousinweb}
\textbf{\bn{লাড্ডু}} (\emph{lāḍḍu}) matches Sanskrit \emph{laḍḍu}, \textquotedblleft{}a kind of sweetmeat\textquotedblright{}; the Bengali word inherited from it is \emph{lāṛu}, which CDIAL lists (CDIAL 10926). CDIAL takes the Sanskrit word as probably the same as \emph{*laḍḍu}, \textquotedblleft{}lump\textquotedblright{}, and notes a proposed link with Munda words for \textquotedblleft{}lump\textquotedblright{}.
\end{cousinweb}

\begin{grammarlens}[title={What you've built}]
One more word for this chapter.
\end{grammarlens}

\subsection*{Your turn}
\begin{itemize}
\raggedright
  \item \emph{Say it:} \emph{lāḍḍu}
  \item \emph{Say it:} \emph{lāḍḍu}, once more
  \item \emph{Say it:} say \emph{lāḍḍu}
  \item \emph{Recall:} say the Bengali for a scholarship, then the Bengali for a jalebi, then say \emph{lāḍḍu} again
\end{itemize}

\begin{tcolorbox}[breakable,colback=teal!4,colframe=teal!35!black,title={Before you move on}]
What does \bn{লাড্ডু} mean? (\textquotedblleft{}A laddu\textquotedblright{}.) Say it once more.
\end{tcolorbox}
\section[\texorpdfstring{\bn{গুড়} (guṛ)}{গুড় (guṛ)}]{\bn{গুড়} (guṛ) — jaggery}

\label{lesson:BN-C173-gur}



\begin{quote}
Before the new one: say the Bengali for a jalebi, then the Bengali for a laddu.
\end{quote}

\subsection*{You'll want to know: \bn{গুড়}}
\textbf{\bn{গুড়}} — \emph{guṛ} — \textquotedblleft{}jaggery\textquotedblright{}.

\begin{cousinweb}
\textbf{\bn{গুড়}} (\emph{guṛ}) is a tadbhava, a word inherited from Sanskrit and reshaped by sound change. It comes from \emph{guḍa}, \textquotedblleft{}boiled sugarcane juice, molasses\textquotedblright{} (CDIAL 4182). CDIAL lists Bengali \emph{guṛ}, \textquotedblleft{}molasses\textquotedblright{}, and asks whether the word is the same as \emph{guḍa}, \textquotedblleft{}ball\textquotedblright{}.
\end{cousinweb}

\begin{grammarlens}[title={What you've built}]
One more word for this chapter.
\end{grammarlens}

\subsection*{Your turn}
\begin{itemize}
\raggedright
  \item \emph{Say it:} \emph{guṛ}
  \item \emph{Say it:} \emph{guṛ}, once more
  \item \emph{Say it:} say \emph{guṛ}
  \item \emph{Recall:} say the Bengali for a jalebi, then the Bengali for a laddu, then say \emph{guṛ} again
\end{itemize}

\begin{tcolorbox}[breakable,colback=teal!4,colframe=teal!35!black,title={Before you move on}]
What does \bn{গুড়} mean? (\textquotedblleft{}Jaggery\textquotedblright{}.) Say it once more.
\end{tcolorbox}
\section[\texorpdfstring{\bn{জাম} (jām)}{জাম (jām)}]{\bn{জাম} (jām) — a black plum}

\label{lesson:BN-C173-jam}



\begin{quote}
Before the new one: say the Bengali for a laddu, then the Bengali for jaggery.
\end{quote}

\subsection*{You'll want to know: \bn{জাম}}
\textbf{\bn{জাম}} — \emph{jām} — \textquotedblleft{}a black plum\textquotedblright{}.

\begin{cousinweb}
\textbf{\bn{জাম}} (\emph{jām}) is a tadbhava, a word inherited from Sanskrit and reshaped by sound change. It comes from \emph{jambū}, \textquotedblleft{}the rose-apple tree\textquotedblright{} (CDIAL 5131). CDIAL lists Bengali \emph{jām} for the tree and its fruit.
\end{cousinweb}

\begin{grammarlens}[title={What you've built}]
One more word for this chapter.
\end{grammarlens}

\subsection*{Your turn}
\begin{itemize}
\raggedright
  \item \emph{Say it:} \emph{jām}
  \item \emph{Say it:} \emph{jām}, once more
  \item \emph{Say it:} say \emph{jām}
  \item \emph{Recall:} say the Bengali for a laddu, then the Bengali for jaggery, then say \emph{jām} again
\end{itemize}

\begin{tcolorbox}[breakable,colback=teal!4,colframe=teal!35!black,title={Before you move on}]
What does \bn{জাম} mean? (\textquotedblleft{}A black plum\textquotedblright{}.) Say it once more.
\end{tcolorbox}
\section[\texorpdfstring{\bn{লিচু} (lichu)}{লিচু (lichu)}]{\bn{লিচু} (lichu) — a lychee}

\label{lesson:BN-C173-lichu}



\begin{quote}
Before the new one: say the Bengali for jaggery, then the Bengali for a black plum.
\end{quote}

\subsection*{You'll want to know: \bn{লিচু}}
\textbf{\bn{লিচু}} — \emph{lichu} — \textquotedblleft{}a lychee\textquotedblright{}.

\begin{cousinweb}
\textbf{\bn{লিচু}} (\emph{lichu}) is usually traced to Chinese \emph{lizhi}, the name of the fruit.
\end{cousinweb}

\begin{grammarlens}[title={What you've built}]
One more word for this chapter.
\end{grammarlens}

\subsection*{Your turn}
\begin{itemize}
\raggedright
  \item \emph{Say it:} \emph{lichu}
  \item \emph{Say it:} \emph{lichu}, once more
  \item \emph{Say it:} say \emph{lichu}
  \item \emph{Recall:} say the Bengali for jaggery, then the Bengali for a black plum, then say \emph{lichu} again
\end{itemize}

\begin{tcolorbox}[breakable,colback=teal!4,colframe=teal!35!black,title={Before you move on}]
What does \bn{লিচু} mean? (\textquotedblleft{}A lychee\textquotedblright{}.) Say it once more.
\end{tcolorbox}
